\documentclass[11pt]{article}
\usepackage[margin=1in]{geometry}
\usepackage{booktabs}
\usepackage{array}
\usepackage{longtable}
\usepackage{hyperref}
\usepackage[T1]{fontenc}
\usepackage[utf8]{inputenc}
\usepackage{lmodern}
\title{Qwen3 Loop-Trigger Attention Pilot}
\author{}
\date{2026-04-13}

\begin{document}
\maketitle

\section*{Question}

This note answers one bounded question on the saved March 22--23 Qwen3-1.7B instruct prompt-profile archives: when the model reaches the final repeated $30$-gram that trips the $n=30$, $k=20$ loop detector, is the trigger token attending to earlier loop copies?

All numbers below come from a corrected rerun after fixing a per-layer wrapper capture bug in the first draft of the attention script.

\section*{Commit Audit}

I fetched \texttt{upstream/main} on 2026-04-13 17:24 UTC. The only new upstream change was merge commit \texttt{d4255f2} (PR \#9, unified prompt-profile report surface). No new loop-trigger attention-analysis script landed there, so the analysis below is a fresh implementation against the saved archives.

\section*{Reconstruction Boundary}

The saved prompt-profile archives preserve prompt token IDs, completion text, finish reason, saved completion length, and saved first-loop prefix length, but not exact completion token IDs. The reconstruction pass over all five saved datasets found:

\begin{center}
\begin{tabular}{lrrrr}
\toprule
Dataset & Rollouts & Loop rows & Exact retokenized length & Trigger-prefix match \\
\midrule
GPQA & 792 & 124 & 46 / 792 & 121 / 124 \\
AIME & 240 & 31 & 38 / 240 & 30 / 31 \\
MATH-500 & 2000 & 68 & 38 / 2000 & 68 / 68 \\
MMLU-Pro & 3200 & 130 & 34 / 3200 & 130 / 130 \\
LiveCodeBench & 3200 & 467 & 356 / 3200 & 462 / 467 \\
\bottomrule
\end{tabular}
\end{center}

The dominant mismatch is exactly one missing hidden stop token on \texttt{finish\_reason=stop}. Qwen3's configured EOS IDs are \texttt{151645 = <|im\_end|>} and \texttt{151643 = <|endoftext|>}; both decode to empty text. So exact completion-ID replay is still not available from the saved text alone, but the trigger prefix is recovered on essentially every loop row.

\section*{Attention Pilot}

I selected the shortest exact-trigger rows with total prefix length at most \texttt{4096}:

\begin{itemize}
  \item GPQA: 3 rows (712, 1598, 1684)
  \item AIME: 2 rows (3686, 3870)
  \item MATH-500: 3 rows (409, 805, 1348)
  \item MMLU-Pro: 3 rows (2034, 2048, 2112)
  \item LiveCodeBench: 3 rows (1544, 1662, 2554)
\end{itemize}

For each row, the probe looks only at the final token of the triggering $30$-gram copy and measures how much attention mass falls on:

\begin{itemize}
  \item all earlier copies of the same loop $30$-gram,
  \item the prompt,
  \item the current trigger span itself,
  \item and the top-1 attention destination of each head.
\end{itemize}

The cleanest read is the final layer:

\begin{center}
\begin{tabular}{lrrrr}
\toprule
Dataset & Rows & Prev-loop mass & Prompt mass & Current-trigger mass \\
\midrule
GPQA & 3 & 0.038 & 0.755 & 0.154 \\
AIME & 2 & 0.035 & 0.724 & 0.170 \\
MATH-500 & 3 & 0.102 & 0.654 & 0.192 \\
MMLU-Pro & 3 & 0.048 & 0.657 & 0.203 \\
LiveCodeBench & 3 & 0.112 & 0.625 & 0.224 \\
\midrule
Overall & 14 & 0.069 & 0.680 & 0.190 \\
\bottomrule
\end{tabular}
\end{center}

\noindent Head destinations make the same point even more strongly: across all 14 rows, 85.3\% of final-layer heads put their top-1 attention destination on the prompt, 14.7\% on the current trigger/self, and 0\% on earlier loop copies.

\medskip

That is not the whole story, though. Earlier loop copies do show up in intermediate layers:

\begin{itemize}
  \item GPQA: best previous-loop layer is layer 1 with mean previous-loop mass 0.275.
  \item AIME: best previous-loop layer is layer 0 with mean previous-loop mass 0.282.
  \item MATH-500: best previous-loop layer is layer 18 with mean previous-loop mass 0.410.
  \item MMLU-Pro: best previous-loop layer is layer 11 with mean previous-loop mass 0.229.
  \item LiveCodeBench: best previous-loop layer is layer 13 with mean previous-loop mass 0.448.
\end{itemize}

\section*{Read}

The corrected answer is more nuanced than the first draft.

\begin{itemize}
  \item If the claim is ``the model completely ignores earlier loop copies,'' that is too strong across the whole stack; several datasets show meaningful previous-loop attention in intermediate layers.
  \item But on this short exact-trigger slice, the final layer is mostly prompt-focused rather than previous-loop-focused.
  \item Previous-loop mass in the final layer is still nonzero, especially on LiveCodeBench and MATH-500, but it is smaller than prompt mass on every dataset in this bounded pilot.
\end{itemize}

So the honest current answer is:

\begin{itemize}
  \item No, the full-stack claim ``it is not paying attention to previous loops'' is too strong.
  \item Yes, on this short exact-trigger slice the final layer is overwhelmingly prompt-focused rather than previous-loop-focused.
  \item The current bounded read is that earlier loop copies are present in the computation, but they are not the dominant late-layer attention target at the trigger.
\end{itemize}

\section*{Example Rows}

\begin{itemize}
  \item GPQA, sample 36, rollout 3, total prefix 712:
  final layer: prev-loop mass 0.030, prompt mass 0.848, current-trigger mass 0.099. Best previous-loop layer is layer 1, where previous-loop mass rises to 0.424 against prompt mass 0.114.
  \item MATH-500, sample 437, rollout 2, total prefix 409:
  final layer: prev-loop mass 0.104, prompt mass 0.664, current-trigger mass 0.167. Best previous-loop layer is layer 16, where previous-loop mass rises to 0.667.
  \item LiveCodeBench, sample 132, rollout 0, total prefix 1544:
  final layer: prev-loop mass 0.077, prompt mass 0.662, current-trigger mass 0.224. Best previous-loop layer is layer 13, where previous-loop mass rises to 0.488 against prompt mass 0.111.
\end{itemize}

\section*{Code Changes For Future Rollouts}

This task also patched both archive paths so future runs stop losing the exact objects this analysis needs.

\begin{itemize}
  \item \texttt{scripts/build\_probe\_dataset.py} now writes \texttt{completion\_token\_ids} and a structured \texttt{loop\_trigger} object into each prompt-profile rollout row.
  \item \texttt{scripts/collect\_model\_stats.py} now writes a compressed sidecar
  \texttt{\_\_rollout\_archive.jsonl.gz} beside the aggregate stats JSON, with prompt text, prompt token IDs, completion text, completion token IDs, finish reason, and structured trigger metadata.
  \item Two review-caught implementation bugs were fixed before publish: resume-from-LiveCodeBench checkpoint now preserves the matching rollout sidecar, and the collector now spills prompt-rollout chunks batch-by-batch instead of buffering the whole archive in memory.
\end{itemize}

\section*{Next Honest Step}

The next step is to repeat the same measurement on longer trigger prefixes once GPU time is available and to add a matched non-loop control slice, but future runs should use the newly saved exact completion token IDs rather than reconstructed text.

\end{document}
